\documentclass[10pt,
% a4paper,
%twocolumn,
fleqn,
%landscape, 
% papersize,
dvipdfmx,
uplatex
]{jsarticle}



\def\maru#1{\textcircled{\scriptsize#1}}%丸囲み番号

% \RequirePackage[2020/09/30]{platexrelease}

%太字設定
\usepackage[deluxe]{otf}

\usepackage{emathEy}

\usepackage[g]{esvect}

%大きな文字
\usepackage{fix-cm}

%定理環境
\usepackage{emathThm}
%\theoremstyle{boxed}
\theorembodyfont{\normalfont}
\newtheorem{Question}{問題}[subsection]
\newtheorem{Q}{}[subsection]
\newtheorem{question}[Question]{}
\newtheorem{quuestion}{}[subsection]

%セクション，大問番号のデザイン
\renewcommand{\labelenumi}{(\arabic{enumi})}
\renewcommand{\theenumii}{\alph{enumii})}
\renewcommand{\thesection}{第\arabic{section}章}

%用紙サイズの詳細設定
\usepackage{bxpapersize}
\papersizesetup{size={80mm,45mm}}
\usepackage[top=0.7zw,bottom=0truemm,left=3truemm,right=133truemm]{geometry}
\usepackage[dvipdfmx]{graphicx}

%余白など
\usepackage{setspace} % 行間
\setlength{\mathindent}{1zw}
\setlength\parindent{0pt}


%色カラーに関する設定
\usepackage{color}
\definecolor{shiro}{rgb}{0.95703125,0.87109375,0.7421875}
\definecolor{kin}{rgb}{0.95703125,0.87109375,0.7421875}
\definecolor{orange}{rgb}{1,0.7,0.2}
\definecolor{bradorange}{rgb}{1,0.5,0}
\definecolor{pink}{rgb}{0.9176,0.5686,0.5960}
\definecolor{mizu}{rgb}{0.6156,0.8,0.9955}
\color{kin}
% \pagecolor{hukamido}

\usepackage{at}%図の配置
% \usepackage{wallpaper}

\begin{document}



\at(0cm,0cm){\includegraphics[width=8cm,bb=0 0 1920 1080]{./youtube/thumbnails/templates/smart_background/場合の数.jpeg}}
{\color{orange}\bf\boldmath\huge\underline{区別する$\cdot$しない}}\vspace{0.3zw}

\scriptsize 
\bf\boldmath 問．$6$個のボールを$3$つの箱に入れるとき，入れ方は何通りか．\vspace{0.1zw}\\
\hfill $\maru{1}$空箱があってもよい$\maru{2}$空箱はなし\ で，それぞれ求めよ．\vspace{0.1zw}\\
(1)  $1$から$6$まで異なる番号のついた$6$個のボールを\vspace{0.1zw}\\
\hfill $\text{A},\;\text{B},\;\text{C}$と区別された$3$つの箱に入れる場合．\vspace{0.1zw}\\
(2)  互いに区別の付かない$6$個のボールを\vspace{0.1zw}\\
\hfill $\text{A},\;\text{B},\;\text{C}$と区別された$3$つの箱に入れる場合．\vspace{0.1zw}\\
(3)  $1$から$6$まで異なる番号のついた$6$個のボールを\vspace{0.1zw}\\
\hfill 区別のつかない$3$つの箱に入れる場合．\vspace{0.1zw}\\
(4)  互いに区別の付かない$6$個のボールを\vspace{0.1zw}\\
\hfill 区別のつかない$3$つの箱に入れる場合．\vspace{0.1zw}\\

\at(7.0cm,0.2cm){\small\color{bradorange}$\overset{\text{場合の数}}{\text{応用}}$}


\newpage



\at(0cm,0cm){\includegraphics[width=8cm,bb=0 0 1920 1080]{./youtube/thumbnails/templates/smart_background/場合の数.jpeg}}
{\color{orange}\bf\boldmath\huge\underline{完全順列}}\vspace{0.3zw}

\small
\bf\boldmath 問．次の人数で

\LARGE
\vspace{-0.2zw}
プレゼント交換するとき，\\
\hfill 受け取り方は何通りあるか．
\vspace{0.2zw}

\small
ただし，全員が他人のプレゼントを受け取るとする．

\small
(1)  $1$人\hfill
(2)  $2$人\hfill
(3)  $3$人\hfill
(4)  $4$人\hfill
(5)  $5$人\hfill
(6)  $6$人\hfill\\

\at(7.0cm,0.2cm){\small\color{bradorange}$\overset{\text{場合の数}}{\text{応用}}$}


\newpage



\at(0cm,0cm){\includegraphics[width=8cm,bb=0 0 1920 1080]{./youtube/thumbnails/templates/smart_background/場合の数.jpeg}}
{\color{orange}\bf\boldmath\huge\underline{重複組合せ Lv.3 }}\vspace{0.3zw}

\scriptsize
\bf\boldmath 問．次の等式$\cdot$不等式を満たす整数の組$\left(x,\;y,\;z\right)$の個数を求めよ．

\normalsize
\vspace{-0.2zw}
(1)  $x+y+z=6,\;x\geqq 0,\;y\geqq 0,\;z\geqq 0$\\
(2)  $x+y+z=6,\;x\geqq 1,\;y\geqq 1,\;z\geqq 1$\\
(3)  $x+y+z\leqq 6,\;x\geqq 0,\;y\geqq 0,\;z\geqq 0$\\
(4)  $1\leqq x<y<z\leqq 6$\\
(5)  $1\leqq x\leqq y\leqq z\leqq 6$\\

\at(7.0cm,0.2cm){\small\color{bradorange}$\overset{\text{場合の数}}{\text{応用}}$}


\newpage



\at(0cm,0cm){\includegraphics[width=8cm,bb=0 0 1920 1080]{./youtube/thumbnails/templates/smart_background/場合の数.jpeg}}
{\color{orange}\bf\boldmath\huge\underline{最短経路}}\vspace{0.3zw}

\LARGE 
\bf\boldmath 問．以下図で$A$地点から$B$地点まで行く最短経路の総数を求めよ．
\at(7.0cm,0.2cm){\small\color{bradorange}$\overset{\text{場合の数}}{\text{応用}}$}


\newpage



\at(0cm,0cm){\includegraphics[width=8cm,bb=0 0 1920 1080]{./youtube/thumbnails/templates/smart_background/場合の数.jpeg}}
{\color{orange}\bf\boldmath\Large\underline{同じものを含む円/数珠順列}}\vspace{0.3zw}

\normalsize
\bf\boldmath 問．白玉$1$個，赤玉$2$個，黄玉$4$個がある．\\
\hspace{-0.2zw}(1)\hspace{-0.2zw}  これらを机の上に円形に並べる方法は何通りか．\hspace{-0.5zw}


\hspace{-0.1zw}(2)

\LARGE
\vspace{-1zw}
\hspace{0.8zw} これらの玉で\\
\hfill 何通りの首飾りができるか．\\

\at(7.0cm,0.2cm){\small\color{bradorange}$\overset{\text{場合の数}}{\text{応用}}$}


\newpage



\at(0cm,0cm){\includegraphics[width=8cm,bb=0 0 1920 1080]{./youtube/thumbnails/templates/smart_background/場合の数.jpeg}}
{\color{orange}\bf\boldmath\huge\underline{立方体の色塗り}}\vspace{0.3zw}

\small 
\bf\boldmath 問．立方体に色を塗る塗り方は全部で何通りあるか求めよ．ただし，隣接する面は異なる色であり，かつ回転したり倒したりして同じになる塗り方は$1$通りとする．

\large
\hspace{-0.1zw}(1)\hspace{-0.2zw}  各面に異なる$6$色をすべて用いて塗る．\\
\hspace{-0.1zw}(2)\hspace{-0.2zw}   各面に異なる$5$色をすべて用いて塗る．\\
\hspace{-0.1zw}(3)\hspace{-0.2zw}   各面に異なる$4$色をすべて用いて塗る．\\

\at(7.0cm,0.2cm){\small\color{bradorange}$\overset{\text{場合の数}}{\text{応用}}$}


\newpage



\at(0cm,0cm){\includegraphics[width=8cm,bb=0 0 1920 1080]{./youtube/thumbnails/templates/smart_background/場合の数.jpeg}}
{\color{orange}\bf\boldmath\Large\underline{整数をつくる問題 Lv.2 }}\vspace{0.3zw}

\normalsize 
\bf\boldmath 問．$9$個の数字

\LARGE 
\vspace{-0.5zw}
\hspace{0.1zw}
$2,\;2,\;2,\;2,\;3,\;3,\;3,\;4,\;4$

\normalsize 
のうち$4$個を使って$4$桁の数をつくる．

\Large
(1)  全部で何個できるか．\\
(2)  $3$の倍数は何個できるか．\\

\at(7.0cm,0.2cm){\small\color{bradorange}$\overset{\text{場合の数}}{\text{応用}}$}


\newpage



\at(0cm,0cm){\includegraphics[width=8cm,bb=0 0 1920 1080]{./youtube/thumbnails/templates/smart_background/場合の数.jpeg}}
{\color{orange}\bf\boldmath\Large\underline{整数をつくる問題 Lv.1 }}\vspace{0.0zw}

\normalsize 
\bf\boldmath 問．

\Huge 
\vspace{-1zw}
\hspace{0.5zw}
$0,\;1,\;2,\;3,\;4,\;5$

\normalsize 
から異なる$3$つの数字を選んで$3$桁の整数を作る．次を求めよ．

\large
(1)  異なる整数の個数\hspace{0.1zw}
(2)  偶数の個数\\
(3)  $3$の倍数の個数\hspace{0.8zw}
(4)  異なる数の総和\\

\at(7.0cm,0.2cm){\small\color{bradorange}$\overset{\text{場合の数}}{\text{応用}}$}


\newpage



\at(0cm,0cm){\includegraphics[width=8cm,bb=0 0 1920 1080]{./youtube/thumbnails/templates/smart_background/場合の数.jpeg}}
{\color{orange}\bf\boldmath\huge\underline{辞書式に並べる}}\vspace{0.3zw}

\Large 
\bf\boldmath 問．a,\;i,\;k,\;o,\;s,\;yの$6$文字を辞書式に一列に並べて，文字列を作る．

\Large
\vspace{0.2zw}
\huge
(1) aoiskyは何番目か．\\
\Large
(2)  ${352}$番目の文字列を求めよ．\\

\at(7.0cm,0.2cm){\small\color{bradorange}$\overset{\text{場合の数}}{\text{応用}}$}


\newpage



\at(0cm,0cm){\includegraphics[width=8cm,bb=0 0 1920 1080]{./youtube/thumbnails/templates/smart_background/場合の数.jpeg}}
{\color{orange}\bf\boldmath\huge\underline{正の約数の個数}}\vspace{0.3zw}

\small 
\bf\boldmath 
(1)  ${5400}$の正の約数の個数と約数の総和を求めよ．\\
(2)  ${10}!$の正の約数の個数を求めよ．\\
(3)  ${30}!$は最後にいくつ$0$が並ぶか．

\Large
(4)  $p$を素数，$n$を正の整数とする．

\LARGE
\vspace{-0.2zw}
\hspace{1zw}$p^n!$は$p$で何回割れるか．\\

\at(7.0cm,0.2cm){\small\color{bradorange}$\overset{\text{場合の数}}{\text{応用}}$}


\end{document}

